\chapter{Notations}\label{sec:notations}

\Xiaohao{To be revised.}
We use $\lesssim$ to denote inequality modulo a generic positive constant, while we use $C$ to denote some generic positive constant whose definition can change from line to line. We will fix some arbitrary $T\in (0,\infty)$ and consider all processes on the time interval $[0,T]$. We denote the $d$-dimensional  torus of unit length by $\mathbb{T}^d = (\mathbb{R}/\mathbb{Z})^d$, and we write $\mathbb{Z}^d_{+}$ for the nonzero elements in $\mathbb{Z}^d$ whose first nonzero entry is positive. We call a function from $\mathbb R^d$ to $\mathbb R$ a kernel on $\mathbb{R}^d$ if it is smooth, non-negative, radially symmetric, supported in the unit ball of $\mathbb{R}^d$ and it has finite moments up to the second order. We use $W^{k,p}$ to denote $L^p$ Sobolev spaces on $\mathbb T^d$. For $p=2$ we use the Fourier definition of Sobolev spaces of any regularity (also fractional or negative), and in particular the inner product on $W^{-1,2}(\mathbb{T}^d)$ is
\begin{equation}\label{eq:inner-prod-H-1}
    \langle u,v \rangle_{W^{-1,2}} = \langle (1-\Delta)^{-1/2}u, (1-\Delta)^{-1/2}v\rangle_{L^2}.
\end{equation}
We write $e_k(x) := e^{2\pi i k \cdot x}$ with $k \in \mathbb{Z}^d$ for the Fourier basis of $L^2(\mathbb{T}^d,\mathbb{C})$. The Fourier transform on $\mathbb R^d$ is defined with the normalization $\hat \varphi (k) = \mathcal{F}_{\mathbb R^d}\varphi(k) = \int_{\mathbb R^d} e^{-2\pi i k \cdot x} \varphi(x) \mathrm{d} x$\\ 

%while $(B_t)_{t \geq 0}$ denotes a $d$-dimensional complex-valued Brownian motion, i.e. with independent components $B(1)\dots, B(d)$ and such that $\operatorname{Re}(B^i)$ and $\operatorname{Im}(B^i)$ are independent real-valued Brownian motions. \\

We will use calculus on the space of probability measures, which slightly differs from \cite{KLvR19} who used calculus on the Radon measures with finite total variation. We denote the space of probability measure on the torus (with Borel $\sigma$-algebra) as $\mcl{P} := \mcl{P}(\mathbb{T}^d)$, and we endow $\mcl{P}$ with the 2-Wasserstein distance
\begin{equation*}
    \mcl{W}_{2}^2(\mu,\nu) := \inf_{\Pi \in \Pi(\mu,\nu)} \int_{\mathbb{T}^d \times \mathbb{T}^d}\mathrm{d}^2_{\mathbb{T}^d}(x,y)\mathrm{d}\Pi(x,y),
\end{equation*}
where $\mathrm{d}_{\mathbb{T}^d}$ denotes the canonical distance on $\mathbb{T}^d$ and $\Pi(\mu,\nu)$ denotes the family of couplings of $\mu$ and $\nu$. Since $d_{\mathbb{T}^d}$ is bounded, this means that we are simply considering the weak topology on $\mcl{P}$. We use $\rightharpoonup$ to indicate weak convergence on the space of probability measures. Functionals on $\mcl{P}$ of the form $\mcl{F} = f\big(\langle \varphi, \cdot \rangle\big)$ are called cylindrical if there exists some $m \in \mathbb{N}$ such that $f: \mathbb{R}^m \to \mathbb{R}$ and $\varphi: \mathbb{T}^d \to \mathbb{R}^m$ are both smooth, and $f$ has bounded derivatives of all orders. For differentiation, $\delta/\delta \mu$ is used to denote the extrinsic derivative of functionals on $\mcl{P}$, while $D_{\mu}$ is used to denote the interior derivative (also called Lions derivative).  \\

We write $|b|^2 := \sum_{i = 1}^d |b_i|^2$ for the $\ell^2$ norm on $\mathbb R^d$. For $d \times d$ matrices $A$ and $B$, we denote by $A^{\transpose}$ the transpose of $A$ and by $A:B := \sum_{i,j=1}^d A_{ij}B_{ij}$ the Frobenius product of $A$ and $B$. We denote the Frobenius norm of $d \times d$ matrix $A$ as $|A| := (\sum_{i,j = 1}^d |A_{ij}|^2)^{1/2}$, for which we have 
\begin{equation}\label{eq:matrix_bound}
    |A:B| \leq |A|\cdot|B|, \quad |A\cdot B| \leq |A|\cdot |B|, \quad |Ab| \leq |A|\cdot|b|,
\end{equation}
all following from the Cauchy-Schwarz inequality. If $X$ is a normed vector space, we extend the definition to $X^{d \times d}$ by setting $\|F\|_{X} := \big|(\|F_{ij}\|_{X})\big|$. We use $L_2$ to denote the Hilbert-Schmidt norm. Given an SPDE in a Gelfand triple $V \subset H \subset V^{\ast}$ in the sense of \cite{LR15}, we define the \emph{$H$ solution} to the SPDE to be the probabilistically strong solution in $C([0,T],H) \cap L^2([0,T],V)$ for any $0< T < \infty$. Throughout the paper, all stochastic integrals are considered in the It\^o sense.


\section*{Function Spaces}
We denote by $\mathbb{T}^d = (\mathbb{R}/\mathbb{Z})^d$ the $d$-dimensional
torus and write $L_{\mathbb{T}^d}^p = L^p (\mathbb{T}^d),$ $p \in [1, \infty]$
for the usual Lebesgue spaces. By $H_{\mathbb{T}^d}^{\alpha}$ we mean the
Sobolev space with index $\alpha \in \mathbb{R}$ defined as
\[
H_{\mathbb{T}^d}^{\alpha} := \left\{ u \in \mathcal{S}'_{\mathbb{T}^d} :
\| (1 - \Delta)^{\alpha / 2} u \|_{L_{\mathbb{T}^3}^2} < \infty \right\}.
\]
Moreover, we write $W_{\mathbb{T}^d}^{s, p} := B_{p, p}^s (\mathbb{T}^d)$
which we call fractional Sobolev spaces and
$\mathcal{C}_{\mathbb{T}^d}^s := B_{\infty, \infty}^s (\mathbb{T}^d)$
called H{\"o}lder-Besov spaces for $s \in \mathbb{R}$, $p \in [1,
\infty]$ which are special cases of Besov spaces whose definition is recalled
in Appendix \ref{sec:append-analy}. Occasionally we will write, by a slight
abuse, $W_{\mathbb{T}^d}^{0, p} = L_{\mathbb{T}^d}^p$ and
$W_{\mathbb{T}^d}^{s, \infty} =\mathcal{C}_{\mathbb{T}^d}^s .$ We also sometimes use the following notation 
\[
\| f (t) \|^p_{L_{t ; [a, b]}^p}
:= \| f \|^p_{L_{[a, b]}^p} := \int^b_a | f (t) |^p d t
\]
for an interval $- \infty \leqslant a < b \leqslant
\infty$ and the obvious modification for $p = \infty$. We recall the definition of Littlewood-Paley blocks and denoted by
$\chi$ and $\rho$ two non-negative smooth and compactly supported radial
functions $\mathbb{R}^d \rightarrow \mathbb{C}$ such that
\begin{enumerate}
  \item The support of $\chi$is contained in a ball and the support of $\rho$
  is contained in an annulus $\{ x \in \mathbb{R}^d : a \leqslant | x |
  \leqslant b \}$
  
  \item For all $\xi \in \mathbb{R}^d$, $\chi (\xi) + \underset{j \geqslant
  0}{\sum} \rho (2^{- j} \xi) = 1$;
  
  \item For $j \geqslant 1$, $\chi (\cdot) \rho (2^{- j} \cdot) = 0$ and
  $\rho (2^{- j} \cdot) \rho (2^{- i} \cdot) = 0$ for $| i - j | > 1$.
\end{enumerate}
The Littlewood-Paley blocks $(\Delta_j)_{j \geqslant - 1}$ associated to $f
\in \mathcal{S}' (\mathbb{T}^d)$ are defined by $\Delta_{- 1} f :=
\mathcal{F}^{- 1} (\chi \mathcal{F} f)$ and $\Delta_j f := \mathcal{F}^{-
1} (\rho (2^{- j} \cdot) \mathcal{F} f)$ for $j \geqslant 0$. We also set,
for $f \in \mathcal{S}' (\mathbb{T}^d)$ and $j \geqslant - 1$ that $S_j f
:= \sum^{j - 1}_{i = - 1} \Delta_i f$. Then the \ Besov space with
parameters $p, q \in [1, \infty), \alpha \in \mathbb{R}$ can now be defined as
$B_{p, q}^{\alpha} (\mathbb{T}^d) := \{ u \in \mathcal{S}' (\mathbb{T}^d)
: \| u \|_{B_{p, q}^{\alpha}} < \infty \}$, where the norm is defined as
\[ \| u \|_{B_{p, q}^{\alpha}} := \left( \underset{k \geqslant - 1}{\sum}
   ((2^{\alpha k} \| \Delta_k u \|_{L^p})^q) \right)^{\frac{1}{q}}, \]
with the obvious modification for $q = \infty .$ We also define the
Besov-H{\"o}lder spaces $\mathcal{C}^{\alpha} := B_{\infty,
\infty}^{\alpha}$, which for $\alpha \in (0, 1)$ agree with the usual
H{\"o}lder spaces $C^{\alpha} .$ Using this notation, we can formally
decompose the product $f \cdot g$ of two distributions $f$ and $g$ as \
\[ f \cdot g = f \olessthan g + f \odot g + f \ogreaterthan g, \]
where
\[ f \olessthan g := \underset{j \geqslant - 1}{\sum} S_{j - 1} f
   \Delta_j g \quad \text{and\quad} f \ogreaterthan g := \underset{j
   \geqslant - 1}{\sum} \Delta_j f S_{j - 1} g \]
are referred to as the paraproducts, whereas
\[ f \odot g := \underset{j \geqslant - 1}{\sum} \underset{| i - j |
   \leqslant 1}{\sum} \Delta_i f \Delta_j g \]
is called the resonant product. An important point is that the
paraproduct terms are always well defined whatever the regularity of $f$ and
$g$. The resonant product, on the other hand, is a priori only well defined if
the sum of their regularities is positive. We collect some results.

\begin{lem}[{\cite{CM17}}, Theorem 3.17]
  \label{lem:para}Let $\alpha, \alpha_1, \alpha_2 \in \mathbb{R}$ and $p, p_1,
  p_2, q \in [1, \infty]$ be such that
  \[ \alpha_1 \neq 0 \quad \alpha = (\alpha_1 \wedge 0) + \alpha_2 \quad
     \text{and} \quad \frac{1}{p} = \frac{1}{p_1} + \frac{1}{p_2} . \]
  Then we have the bound
  \[ \| f \olessthan g \|_{B_{p, q}^{\alpha}} \lesssim \| f \|_{B_{p_1,
     \infty}^{\alpha_1}} \| g \|_{{B^{\alpha_2}_{p_2, q}} } \]
  and in the case where $\alpha_1 + \alpha_2 > 0$ we have the bound
  \[ \| f \odot g \|_{B_{p, q}^{\alpha_1 + \alpha_2}} \lesssim \| f
     \|_{B_{p_1, \infty}^{\alpha_1}} \| g \|_{{B^{\alpha_2}_{p_2, q}} } . \]
\end{lem}

\begin{rem}
  In this paper, we only consider the cases $W^{s, p} {= B_{p, p}^s}
  $ fractional Sobolev, $H^s {= B_{2, 2}^s} $ and
 H{\"o}lder Besov spaces $\mathcal{C}^s = B_{\infty, \infty}^s$
  using the convention $W^{s, 2} = H^s$ and $W^{s, \infty} =\mathcal{C}^s$.
\end{rem}

\begin{lem}[Bernstein Inequality]
  \label{lem:bernstein} Let $\mathcal{A}$ be an annulus and $\mathcal{B}$ be a
  ball. For any $k \in \mathbb{N}, \lambda > 0,$and $1 \leqslant p \leqslant q
  \leqslant \infty$ we have
  \begin{enumerate}
    \item if $u \in L^p (\mathbb{T}^d) $is such that $\mathrm{supp}
    (\mathcal{F}u) \subset \lambda \mathcal{B}$ then
    \[ \underset{\mu \in \mathbb{N}^d : | \mu | = k}{\max} \| \partial^{\mu} u
       \|_{L^q} \lesssim_k \lambda^{k + d \left( \frac{1}{p} - \frac{1}{q}
       \right)} \| u \|_{L^p} \]
    \item if $u \in L^p (\mathbb{T}^d) $is such that $\mathrm{supp}
    (\mathcal{F}u) \subset \lambda \mathcal{A}$ then
    \[ \lambda^k \| u \|_{L^p} \lesssim_k \underset{\mu \in \mathbb{N}^d : |
       \mu | = k}{\max} \| \partial^{\mu} u \|_{L^p} . \]
  \end{enumerate}
\end{lem}

\

\begin{lem}[Besov Embedding]
  \label{lem:besovem} Let $\alpha < \beta \in \mathbb{R}, $ $q_1 \leqslant
  q_2$, and $p > r \in [1, \infty]$ be such that $\beta = \alpha + d (1 / r -
  1 / p)$, then we have the following bound
  \[ \| f \|_{B_{p, q_2}^{\alpha} (\mathbb{T}^d)} \lesssim \| f \|_{{B_{r,
     q}^{\beta}}_1 (\mathbb{T}^d)} . \]
  Moreover, we have the following embedding between fractional Sobolev spaces
  and Bessel potential spaces
  \begin{equation}
    \| f \|_{W_{\mathbb{T}^d}^{\alpha, \rho}} \lesssim \left\| (1 -
    \Delta)^{\frac{\beta}{2}} f \right\|_{L_{\mathbb{T}^d}^{\rho}} \lesssim \|
    f \|_{W_{\mathbb{T}^d}^{\beta, \rho}} \quad \text{for } \alpha < \beta
    \quad \text{and } p \geqslant 2. \label{ineq:Bessel}
  \end{equation}
\end{lem}

\

\begin{prop}[Commutator in {\cite{GIP15}}, see Proposition A.9
{\cite{MW17}}]
  \label{prop:commu}
  
  Given $\alpha \in (0, 1)$, $\beta, \gamma \in \mathbb{R}$ such that $\beta +
  \gamma < 0$ and $\alpha + \beta + \gamma > 0$, the following trilinear
  operator $C$ defined for any smooth functions $f, g, h$ by
  \[ \mathrm{com} (f, g, h) := (f \olessthan g) \odot h - f \cdot (g \odot
     h) \]
  can be extended continuously to the product space $W^{\alpha, p} \times
  \mathcal{C}^{\beta} \times \mathcal{C}^{\gamma}$. Moreover, we have the
  following bound
  \[ || \mathrm{com} (f, g, h) ||_{W^{\alpha + \beta + \gamma - \delta, p}}
     \lesssim || f ||_{W^{\alpha, p}} || g ||_{\mathcal{C}^{\beta}} || h
     ||_{\mathcal{C}^{\gamma}} \]
  for all $f \in W^{\alpha, p}$, $g \in \mathcal{C}^{\beta}$ and $h \in
  \mathcal{C}^{\gamma}$, and every $\delta > 0, p \geqslant 1$.
\end{prop}

\begin{lem}[Fractional Leibniz {\cite{G96}}]
  \label{lem:fracleib} Let $1 < p < \infty$ and $p_1, p_2, p'_1, p'_2$ such
  that
  \[ \frac{1}{p_1} + \frac{1}{p_2} = \frac{1}{p'_1} + \frac{1}{p'_2} =
     \frac{1}{p} . \]
  Then for any $s, \alpha \geqslant 0$ there exists a constant s.t.
  \[ \| \langle \nabla \rangle^s (f g) \|_{L^p} \leqslant C \| \langle \nabla
     \rangle^{s + \alpha} f \|_{L^{p_2}} \| \langle \nabla \rangle^{- \alpha}
     g \|_{L^{p_1}} + C \| \langle \nabla \rangle^{- \alpha} f \|_{L^{p'_2}}
     \| \langle \nabla \rangle^{s + \alpha} g \|_{L^{p'_1}} . \]
\end{lem}

\section*{Density Matrices}

Let $\mathcal{L}^1_n = \mathcal{L}^1 (L_s^2 ((\mathbb{T}^3)^n))$ be the set of
trace-class operators on the symmetric space $L_s^2 ((\mathbb{T}^3)^n)$, the
set of symmetric $L^2$ functions on $n$-copies of the three dimensional torus
$\mathbb{T}^3$ for $n \in \mathbb{N}$. The class of density matrices
is then the subset of $\rho \in \mathcal{L}^1_n$ such that $\rho$ is
self-adjoint, non-negative and normalized by $\mathrm{Tr} (\rho) = 1$. If $S$ is
a linear operator on $L^2 (\mathbb{T}^3)$, we abbreviate $S_{x_k} = S_k$ to be the
linear operator on the $k$-th component of $L_s^2 ((\mathbb{T}^3)^n)$, where
the tensor product is symmetrized. We also denote the family of compact
operators on $L_s^2 ((\mathbb{T}^3)^n)$ as 
\begin{equation}\label{eq:def-K_n}
\mathcal{K}_n = \mathcal{K} (L_s^2
((\mathbb{T}^3)^n)).
\end{equation}
More generally, we denote by $\mathcal{L}^p (L^2_s ((\mathbb{T}^3)^n))$ the
space of symmetric compact operators with finite $p$-Schatten norm,
namely
\[ A \in \mathcal{L}^p (L^2_s ((\mathbb{T}^3)^n)) \Leftrightarrow \| A
   \|_{\mathcal{L}^p}^p = \mathrm{Tr} (| A |^p) < + \infty . \]
If $A \in \mathcal{L}^p (L^2_s ((\mathbb{T}^3)^n))$ for some $p \leqslant 2$,
$A$ is in particular a Hilbert-Schmidt operator and thus, we can - and will -
identify $A$ with its kernel in $L^2 ((\mathbb{T}^3)^{2 n})$. We further
define the partial trace $\mathrm{Tr}_{k + 1 : n}$ as a bounded operator from
$\mathcal{L}^1 (L^2_s ((\mathbb{T}^3)^n))$ into $\mathcal{L}^1 (L^2_s
((\mathbb{T}^3) )^k)$ for $k \leqslant n$ to be such that
\begin{equation}
  \mathrm{Tr} (\mathrm{Tr}_{k + 1 : n} A K) = \mathrm{Tr} (A K \otimes \mathrm{Id}_{k
  + 1 : n}), \label{eq:def-partialtrace-weak}
\end{equation}
for any $K \in \mathcal{K} (L_s^2 ((\mathbb{T}^3)^{n - k}))$. Equivalently,
the partial trace is defined in terms of kernels by
\[ \mathrm{Tr}_{k + 1 : n} (A) (x_{1 : k}, x'_{1 : k}) = \int_{\mathbb{T}^{3 (n
   - k)}} A ((x_{1 : k}, z_{k + 1 : n}), (x'_{1 : k}, z_{k + 1 : n})) \mathrm{d}
   z_{k + 1 : n} \]
where, hereafter, $w_{h : \ell}$ denotes the vector $w_{h : \ell} \in
(\mathbb{T}^3)^{(\ell - h + 1)}$ formed by the components from $h$-th to
$\ell$-th of an arbitrary vector $w_{1 : n} \in (\mathbb{T}^3)^n$. We call a
family of trace-class operators $(\rho^n)_{n \in \mathbb{N}}$ with $\rho^n \in
\mathcal{L}^1 (L^2_s (\mathbb{T}^{3 n}))$ compatible, if for any $n
\in \mathbb{N}$ and $k \leqslant n$ we have
\[ \rho^k = \mathrm{Tr}_{k + 1 : n} (\rho^n) . \]
We use the same name for a finite set $(\rho_n)_{n \leqslant N}$ satisfying
the equation as above. We further define the Sobolev spaces of operators with
respect to the self-adjoint, negative and densely defined $\mathcal{H}$ and
$\mathcal{H}_{\delta}$ to be
\begin{equation}
\begin{split}
  \mathcal{W}_N^{\alpha, \beta} = & \left\{ \rho \in \mathcal{L}^{\beta}
  ((\mathbb{T}^3)^N) : \| \rho \|_{\mathcal{W}^{\alpha, \beta}_N} < + \infty
  \right\}, \\
  & \text{where } {\| \rho \|_{\mathcal{W}_N^{\alpha, \beta}}}  = \left( \sum_{k
  = 1}^N \mathrm{tr} (| (- \mathcal{H}_{x_k})^{\alpha / 2} \rho (-
  \mathcal{H}_{x_k})^{\alpha / 2} |^{\beta}) \right)^{\frac{1}{\beta}},
  \label{operator-norm-sobolev}
  \end{split}
\end{equation}
and respectively 
\begin{equation}
\begin{split}
  \mathcal{W}_{N, \delta}^{\alpha, \beta} = & \left\{ \rho \in \mathcal{L}^{\beta}
  ((\mathbb{T}^3)^N) : \| \rho \|_{\mathcal{W}^{\alpha, \beta}_{N, \delta}} <
  + \infty \right\}, \\
  & \text{where }\| \rho \|_{\mathcal{W}_{N, \delta}^{\alpha, \beta}} =
  \left( \sum_{k = 1}^N \mathrm{tr} (| (- \mathcal{H}_{\delta, x_k})^{\alpha /
  2} \rho (- \mathcal{H}_{\delta, x_k})^{\alpha / 2} |^{\beta})
  \right)^{\frac{1}{\beta}}, \label{operator-norm-sobolev-delta}
  \end{split}
\end{equation}
where the Anderson Hamiltonian $\mathcal{H}$ is given in Proposition
\ref{prop:construction-AH}, and its mollified version $\mathcal{H}_{\delta}$
is defined in (\ref{eq:mollified-AH}).

\section*{Other notational conventions}

We use some usual conventions, such as writing $a_+ := \max \{ a, 0 \}$, and
\[ a \lesssim b \]
if there exists a constant $C > 0$ s.t. $a \leqslant C b$ which does not
depend on relevant parameters. We also write
\[ a \asymp b \text{ if } a \lesssim b \text{ and } b \lesssim a \]
and if we want to keep track that the implicit constant depends on a parameter
$\gamma$ we write $a \lesssim_{\gamma} b$ etc.

Since the problem we are considering is path-wise, meaning that we
fix a realization $\omega \in \Omega$ of the noise which has an enhanced noise
$\Xi (\omega)$ with $\| \Xi \|_{\mathfrak{X}} < \infty $, see Lemma
\ref{lem:conv-enhancement}. Since all our constructions depend on this
realisation so we allow every implicit constant to depend on this
norm, schematically written as
\[ a \lesssim b \Leftrightarrow a \lesssim_{\Xi} b. \]
On a related note, we often write things like $f \in
\mathcal{C}_{\mathbb{T}^d}^{\alpha -}$ to mean $f \in
\mathcal{C}_{\mathbb{T}^d}^{\beta}$ for all $\beta < \alpha$ and in some cases
we take it to mean $f \in \mathcal{C}_{\mathbb{T}^d}^{\alpha - \kappa}$ for an
arbitrarily small $\kappa > 0$ which can be chosen independently of other
parameters. Similarly, we write something of the form
\[ \| A f \|_X \lesssim \| f \|_{H^{\alpha +}_{\mathbb{T}^d}} \Leftrightarrow
   \| A f \|_X \lesssim \| f \|_{H^{\alpha + \varepsilon}_{\mathbb{T}^d}} 
  \ \mathrm{for all} \ \varepsilon > 0 \]
to mean that the bound is true for every $\varepsilon > 0$ and the implicit
constant blows up as $\varepsilon \to 0^+.$

\section*{Flattened, sharpened and natural variables}

We try to stay consistent in labelling the variables by the ``flat'',
``sharp'', and ``natural'' superscript to denote the transformation one is
considering which are succesively more refined and complicated. For the
convenience of the reader, we briefly summarize their definitions and where
they are used.
\begin{itemize}
  \item ``Flat transform'': multiplication by an exponential of a
  noise term, see \eqref{def:flat}, which characterizes the form domain of
  $\mathcal{H}$ i.e. 
  \begin{equation*}
  u^{\flat} = e^{- W} u \ \mathrm{and} \ \| u \|_{\mathcal{D} \left(
  \sqrt{- \mathcal{H}} \right)} \asymp \| u^{\flat} \|_{H_{\mathbb{T}^3}^1}.
  \end{equation*}
  
  \item ``Sharp transform'': paracontrolled ansatz (inverse of
  $\Gamma$, see Lemma \ref{lem:con-Gamma}) performed after the exponential
  ansatz i.e. $u^{\sharp} = \Gamma^{- 1} e^{- W} u$ which characterizes the
  domain of $\mathcal{H}$ which we also transform as $\mathcal{H}^{\sharp}
  := \Gamma^{- 1} e^{- W} \mathcal{H} e^W \Gamma$ via
  \[ \| u \|_{\mathcal{D} (\mathcal{H})} \asymp \| u^{\sharp}
     \|_{H_{\mathbb{T}^3}^2} \asymp \| \mathcal{H}^{\sharp} u^{\sharp}
     \|_{L_{\mathbb{T}^3}^2} . \]
  We use this transformation in Section \ref{sec:LWPenergy} and to obtain the
  embeddings and norm equivalences in Lemma \ref{lem:embed}.
  
  \item ``Natural transform'': The $\Theta$ transform, see Lemma
  \ref{lem:homeom-Theta}, applied after the sharp transform i.e. $u^{\natural}
  = \Theta u^{\sharp} = \Theta \Gamma^{- 1} e^{- W} u$ which is the inverse of
  $\Lambda$ from Lemma \ref{lem:Lambda}. We set also $\mathcal{H}^{\natural}
  := \Lambda^{- 1} \mathcal{H} \Lambda$ and this transformation has the
  important properties
  \[ \mathcal{H}^{\natural} - \Delta : H_{\mathbb{T}^3}^{\frac{3}{2} +}
     \rightarrow H_{\mathbb{T}^3}^{0 +} \quad \text{and } \Lambda - \mathrm{id}
     : H_{\mathbb{T}^3}^{0 -} \rightarrow H_{\mathbb{T}^3}^{\frac{1}{2} -}, \]
  which are properties that are not true without this additional refinement.
  We heavily use this transform in Section \ref{sec:high-order-energy} and
  this transform is one of the novelties of the current paper, whereas the
  other two already appeared in {\cite{GUZ20}}.
\end{itemize}